\documentclass[conference]{IEEEtran}

\usepackage{graphicx}
\usepackage{amsmath,amssymb}
\usepackage{float}
\usepackage{booktabs}
\usepackage{multirow}

\begin{document}

\title{
    \textbf{Benchmark Report: 541.leela\_r} \\
    \text{Register Spill Analysis on RISC-V via gem5}
}

\author{
    \IEEEauthorblockN{Lütfullah Burak Kaya}
    \IEEEauthorblockA{
        Ozyegin University \\
        burak.kaya.50711@ozu.edu.tr
    }
    \and
    \IEEEauthorblockN{Ismail Akturk}
    \IEEEauthorblockA{
        Ozyegin University \\
        ismail.akturk@ozyegin.edu.tr
    }
}

\newcommand{\LongDate}{%
    \number\day\space
    \ifcase\month
    \or January\or February\or March\or April\or May\or June%
    \or July\or August\or September\or October\or November\or December%
    \fi
    \space \number\year
}

\twocolumn[
\begin{@twocolumnfalse}
\maketitle
\begin{center}{\small \LongDate}\end{center}
\vspace{1em}
\end{@twocolumnfalse}
]

% ================================================

\begin{abstract}
This report presents the register spill characterization of the
\texttt{541.leela\_r} benchmark from SPEC CPU2017, executed on a RISC-V
(RV64GC) target using the gem5 TimingSimpleCPU simulator.
A custom SpillDetector was integrated into gem5 to count store-then-load
patterns on the stack within a user-defined Region of Interest (ROI).
The test dataset (\texttt{test.sgf}, 34 lines) was used as the input.
Results show a spill rate of \textbf{4.44\%} within the ROI,
corresponding to 1.16 billion detected spill events across 26 billion
ROI instructions. The benchmark's Monte Carlo Tree Search (MCTS) engine
for the game of Go creates deep, irregular call stacks with high register
pressure at each node expansion.
\end{abstract}

% =========================================================
\section{Benchmark Overview}

\texttt{541.leela\_r} is the SPEC CPU2017 rate variant of Leela, a
computer Go engine based on the Upper Confidence Trees (UCT) algorithm,
a variant of Monte Carlo Tree Search (MCTS). Given a game record in
Smart Game Format (SGF), Leela replays each position and, for each
move, runs a UCT search to find the best response. The engine exercises
recursive tree traversal, board state evaluation, random simulation
(rollout), and backpropagation of scores --- all with irregular control
flow and deep call stacks.

The RISC-V binary (\texttt{leela\_r.riscv}) was compiled with gem5 ROI
markers inserted in \texttt{src/Leela.cpp}, wrapping the main
game-processing loop:

\begin{itemize}
    \item \texttt{m5\_work\_begin(0,0)} --- before the outer SGF-game loop
    \item \texttt{m5\_work\_end(0,0)}   --- after the loop (in the exception handler)
\end{itemize}

For each game in the SGF file, the ROI loop loads the game tree,
replays each move using \texttt{UCTSearch::think()}, and terminates
when the game ends by resignation or two consecutive passes.

% =========================================================
\section{Simulation Setup}

\subsection{Input Datasets}

The three dataset sizes use different SGF game record files with
different numbers of moves (Table~\ref{tab:inputs}). Each SGF game
requires many UCT search iterations per move.

\begin{table}[H]
\centering
\caption{leela\_r Input Dataset Sizes}
\label{tab:inputs}
\begin{tabular}{lcc}
\toprule
\textbf{Dataset} & \textbf{SGF file} & \textbf{File size (lines)} \\
\midrule
test    & test.sgf  &  34 \\
train   & train.sgf &  20 \\
refrate & ref.sgf   & 198 \\
\bottomrule
\end{tabular}
\end{table}

The \textbf{test} dataset (\texttt{test.sgf}, 34 lines) was selected
for this run. The SGF file encodes one or more Go game records;
Leela processes each move in sequence, running a full UCT search tree
for each position.

\subsection{gem5 Configuration}

\begin{itemize}
    \item \textbf{CPU model:}   TimingSimpleCPU
    \item \textbf{ISA:}         RISC-V RV64GC
    \item \textbf{Caches:}      enabled (default L1 I/D + L2)
    \item \textbf{Memory:}      4\,GB simulated physical RAM
    \item \textbf{Spill mode:}  summary only (\texttt{spill\_verbose=False})
\end{itemize}

The full gem5 invocation was:

\begin{verbatim}
./build/RISCV/gem5.opt
  --outdir=benchmarks/leela/m5out_test
  configs/deprecated/example/se.py
  --cmd=benchmarks/leela/leela_r.riscv
  --options="benchmarks/leela/test.sgf"
  --cpu-type=TimingSimpleCPU
  --caches
  --mem-size=4GB
\end{verbatim}

% =========================================================
\section{Simulation Results}

The simulation completed successfully. Leela processed all games in
\texttt{test.sgf}, printing board states and move selections, and
exited cleanly via the expected end-of-file exception.

\subsection{Pre-ROI Context}

Before the ROI, the binary performed engine initialisation: loading
Zobrist hash tables, setting up the pattern matcher
(\texttt{Matcher::get\_Matcher()}), and allocating the main game state.
Table~\ref{tab:preroi} shows the global counters at ROI entry.

\begin{table}[H]
\centering
\caption{Global Counters at ROI Entry}
\label{tab:preroi}
\begin{tabular}{lr}
\toprule
\textbf{Metric} & \textbf{Value} \\
\midrule
Total instructions (pre-ROI) & 929,649,876 \\
Total spills detected (pre-ROI) & 86,033,182 \\
Pre-ROI spill rate & 9.25\% \\
\bottomrule
\end{tabular}
\end{table}

The pre-ROI spill rate of 9.25\% is notably high, reflecting the
register pressure during Leela's pattern matcher initialisation, which
involves processing large lookup tables with complex indexing.

\subsection{ROI Spill Statistics}

Table~\ref{tab:roi} presents the spill statistics collected within
the ROI.

\begin{table}[H]
\centering
\caption{ROI Spill Statistics --- 541.leela\_r (test.sgf)}
\label{tab:roi}
\begin{tabular}{lr}
\toprule
\textbf{Metric} & \textbf{Value} \\
\midrule
ROI instructions         & 26,063,891,645 \\
ROI stores               &  1,756,742,699 \\
ROI loads                &  5,110,627,713 \\
\midrule
\textbf{ROI spills}      & \textbf{1,158,236,271} \\
\textbf{Spill rate}      & \textbf{4.4438\%} \\
\bottomrule
\end{tabular}
\end{table}

\subsection{Timing}

\begin{table}[H]
\centering
\caption{Simulation Timing}
\label{tab:timing}
\begin{tabular}{lr}
\toprule
\textbf{Metric} & \textbf{Value} \\
\midrule
Simulated time           & 45.78\,s \\
Host wall time           & 14,637\,s ($\approx$4.1\,hours) \\
Total simulated instructions & 26,994,306,259 \\
\bottomrule
\end{tabular}
\end{table}

% =========================================================
\section{Analysis}

\subsection{Spill Rate Interpretation}

The measured spill rate of \textbf{4.44\%} means that approximately
1 in every 23 instructions inside the ROI is a register spill.
Table~\ref{tab:compare} places this among all evaluated benchmarks.

\begin{table}[H]
\centering
\caption{Cross-Benchmark Spill Rate Comparison}
\label{tab:compare}
\begin{tabular}{lrr}
\toprule
\textbf{Benchmark} & \textbf{ROI Spills} & \textbf{Spill Rate} \\
\midrule
507.cactuBSSN\_r & 3,782,099,438 & 7.8756\% \\
500.perlbench\_r &     1,070,055 & 7.2838\% \\
526.blender\_r   &    61,845,869 & 6.2193\% \\
541.leela\_r     & 1,158,236,271 & \textbf{4.4438\%} \\
519.lbm\_r       &   274,130,440 & 4.2130\% \\
544.nab\_r       &   151,310,616 & 2.2612\% \\
538.imagick\_r   &     2,424,899 & 2.0601\% \\
505.mcf\_r       &   446,113,136 & 1.8263\% \\
508.namd\_r      &   327,718,955 & 1.0269\% \\
\bottomrule
\end{tabular}
\end{table}

\subsection{Store vs.\ Load Balance}

The load count (5.11B) is approximately $2.91\times$ higher than the
store count (1.76B). This asymmetry is characteristic of tree-search
workloads: during UCT rollouts, the engine reads board positions,
move histories, and node statistics far more often than it writes them.
Each node expansion reads the parent node's statistics for the UCT
selection formula and loads the current board state for move generation,
while only the selected child's statistics are updated.

\subsection{Spill Fraction of Loads}

Within the ROI, the ratio of spills to total loads is:
\[
\frac{1{,}158{,}236{,}271}{5{,}110{,}627{,}713} \approx 22.7\%
\]
Nearly 1 in 4 load instructions is a spill reload. UCT search involves
deeply nested function calls: \texttt{UCTSearch::think()} calls
\texttt{UCTSearch::playout()}, which recursively calls
\texttt{UCTNode::uct\_select()} and ultimately invokes the board
simulation. Each call frame must save and restore multiple live
variables (node pointers, board state references, visit counts,
score accumulators), causing pervasive spilling throughout the search.

\subsection{Pre-ROI Spill Density}

The pre-ROI spill rate of 9.25\% ($86{,}033{,}182$ spills over
$929{,}649{,}876$ instructions) is the second-highest pre-ROI rate
observed across all benchmarks (after cactuBSSN at 11.83\%).
The Leela pattern matcher initialisation, which builds a large
board pattern lookup table using recursive feature hashing, is
itself computationally intensive and register-pressure-heavy.

\subsection{ROI Coverage}

The ROI accounts for $26{,}063{,}891{,}645$ out of
$26{,}994{,}306{,}259$ total simulated instructions:
\[
\frac{26{,}063{,}891{,}645}{26{,}994{,}306{,}259} \approx 96.6\%
\]
The ROI captures 96.6\% of total execution, confirming that the
game-processing loop dominates the benchmark's instruction footprint.

% =========================================================
\section{Conclusion}

The \texttt{541.leela\_r} benchmark exhibits a register spill rate of
4.44\% within its UCT game-processing ROI on RISC-V RV64GC. With 1.16
billion spill events across 26 billion ROI instructions, Leela's
Monte Carlo tree search creates sustained register pressure through
its deep and irregular call structure. The 2.91:1 load-to-store ratio
and 22.7\% spill fraction of loads both reflect the read-dominated
nature of tree traversal. Notably, the pre-ROI initialisation phase
also exhibits a high spill rate (9.25\%), indicating that pattern
matcher construction is itself register-intensive. These results
position \texttt{541.leela\_r} as a mid-range spill benchmark, placing
fourth among all evaluated SPEC CPU2017 workloads.

\end{document}
